\documentclass[11pt,a4paper]{article}
\usepackage[margin=25mm]{geometry}
\usepackage{amsmath,amssymb,amsthm,booktabs}
\newtheorem{theorem}{Theorem}
\newtheorem{corollary}[theorem]{Corollary}
\newcommand{\gen}[1]{\langle #1\rangle}
\title{The degree sixteen boundary}
\author{Human Compression project: a component of the L044 recovery}
\date{}
\begin{document}\maketitle

Put $V=\mathbb F_2^3$, and identify sixteen points with
$V\times\mathbb F_2$. The binary partition argument previously
established for L044 leads to the following permutation group:
\[
 W=\{(a,b)\mapsto(a+u,b+v(a)):u\in V,\quad
                  v:V\to\mathbb F_2,\quad\sum_{a\in V}v(a)=0\}.
                                                        \tag{1}
\]
We resolve its remaining finite obligation. All group operations
in the certificate are exact permutation operations in GAP.
The proof of the reduction to (1) is acquired separately; it is
not replaced by this computation.

\begin{theorem}\label{finite}
No transitive subgroup of $W$ has a string C-marking of rank five.
Disconnected strings, including any inherited five-label subset
of a longer string, are included.
\end{theorem}
\begin{proof}
We give the exhaustive algorithm and its output. A marking
$(x_i)_{i\in I}$ has the intersection property if
\[
 \gen{x_i:i\in J}\cap\gen{x_i:i\in K}
              =\gen{x_i:i\in J\cap K}\qquad(J,K\subseteq I).
                                                        \tag{2}
\]
Thus every generator is a nonidentity involution outside the
subgroup generated by the others. Every subfamily also satisfies (2).
In a string, only consecutive generators can fail to commute.
Separate its noncommuting components, and reorder the components
by decreasing size. Reordering does not change the generated group
or (2), and distinct components commute. Within a component we
retain its linear order. Consequently the possible component sizes
for five generators are
\[
 5,\quad4+1,\quad3+2,\quad3+1+1,\quad2+2+1,\quad
 2+1+1+1,\quad1+1+1+1+1.                         \tag{3}
\]
The last case is impossible: (2) gives an elementary abelian group
of order $32$. A faithful transitive action of an abelian group is
regular, since an element fixing one point fixes all its translates.
It cannot have degree sixteen. The six other cases are enumerated.

There are $2^7$ even functions $v$, and eight translations, so
$|W|=1024$. A nonidentity element with $u=0$ is an involution,
giving $127$ such elements. For each $u\ne0$, the involution
condition is $v(a+u)=v(a)$; the four pairs of the translation allow
$16$ functions, all of even weight. There are therefore
$127+7\cdot16=239$ involutions, precisely the initial search domain.

For orbit reduction take the permutation group $N$ generated by
all eight individual pair flips, the three coordinate translations,
a cyclic permutation of the coordinates, and the transvection
$(a_0,a_1,a_2)\mapsto(a_0+a_1,a_1,a_2)$. Each generator normalizes
$W$. Affine changes of $a$ preserve even weight; conjugation by
a flip function $w$ adds $w(a)+w(a+u)$ to $v$, whose sum is zero.
The supplied generators have order $344064$ as a group, checked
directly in the model. The argument needs only normalization,
not an assertion that $N$ is the full normalizer or automorphism group.

The six label sets used are the rows of the table below. Gaps
separate components. In each row the labels are visited in positions
$1,5,2,4,3$. Suppose the current chosen labels are $S$, their
generated subgroup is $H$, and their pointwise stabilizer under
conjugation by $N$ is
$C=C_N(x_i:i\in S)$. For the next label $j$, form the complete
set of the $239$ involutions $x\in W$ satisfying
\[
 x\notin H,\qquad
 xx_i=x_ix\ \Longleftrightarrow\ |i-j|>1\quad(i\in S).       \tag{4}
\]
This set is invariant under $C$. The program forms its actual
$C$-orbits, checks their disjoint union against the entire candidate
set, and takes one representative from each. It retains a
representative exactly when the enlarged family satisfies (2).
Recursion replaces $H$ by the enlarged subgroup and $C$ by
$C_C(x)$. At depth five it tests actual transitivity on all sixteen
points. No requirement that the subgroup equal $W$ is imposed.

This recursion is exhaustive by induction on the depth. Any
qualifying full family has a retained restriction by (2); its
next member lies in (4). Conjugating by the current stabilizer
sends that member to one of the representatives while leaving
the already chosen family fixed. Such conjugation preserves all
remaining conditions. Hence some complete branch represents
the full family. Neither subgroup order targets nor time caps
are used to infer a completed search.

For every proposed extension, the reference routine explicitly
constructs the subgroups for all subsets of the chosen family.
Their sets of elements are bit vectors in the fixed list of all
$1024$ elements of $W$. For every unordered pair of subsets it
compares their bit-vector intersection with the subgroup for the
intersection of the label subsets. This is exactly (2), including
the empty subset. A second routine uses the acquired recursive
end-intersection criterion on each consecutive component and
checks that their commuting product has order equal to the product
of their orders. The two decisions are compared before every
pruning step. A disagreement aborts the computation. Thus even
rejected extensions are checked against (2); the reference decision
is not restricted to the eventual accepted leaves.

\begin{center}\small
\begin{tabular}{llllr}
\toprule
Sizes & Labels & Nodes at depths $0,\ldots,5$ & Leaves & Transitive\\
\midrule
$5$ & $1,2,3,4,5$ & $1,6,76,76,32,0$ & $0$ & $0$\\
$4+1$ & $1,2,3,4,6$ & $1,6,76,76,528,0$ & $0$ & $0$\\
$3+2$ & $1,2,3,5,6$ & $1,6,76,76,32,0$ & $0$ & $0$\\
$3+1+1$ & $1,2,3,5,7$ & $1,6,76,76,528,5376$ & $5376$ & $0$\\
$2+2+1$ & $1,2,4,5,7$ & $1,6,76,76,528,0$ & $0$ & $0$\\
$2+1+1+1$ & $1,2,4,6,8$ & $1,6,76,76,528,2912$ & $2912$ & $0$\\
\bottomrule
\end{tabular}
\end{center}

All six searches terminate. The reference routine is called
$20478$ times and performs $5413998$ subset-pair comparisons.
None of the $8288$ complete families is transitive. This excludes
all cases in (3), proving the theorem.
\end{proof}

\begin{corollary}
Let $A_n$ in its natural action have a string C-marking of rank
$r\ge n/2$, where $n\ge16$. Suppose every transitive maximal
omission preserves a partition into two blocks. Then every
double omission is intransitive.
\end{corollary}
\begin{proof}
The acquired binary partition theorem proves that a transitive
double omission under these hypotheses forces $n=16$, $r=8$,
and a transitive subgroup of (1) generated by five of the original
distinguished generators. Its full intersection property is
inherited. Theorem~\ref{finite} excludes that subgroup.
\end{proof}

\paragraph{Replay and controls.}
The files \texttt{model.g}, \texttt{enumeration.g}, and
\texttt{run.g} implement the argument above. Running
\texttt{python3 check.py NEW\_DIRECTORY} copies those files to a
fresh directory and performs the full computation. It requires
all completion markers, an empty error stream, the complete table,
and unchanged source hashes. The node and time caps cause failure,
never acceptance. The four commuting generators consisting of
the coordinate translations and the simultaneous flip of all
pairs give a positive control: their subgroup is regular of order
sixteen and satisfies (2). Repeating one translation gives a
negative control which must fail (2). Control comparisons are
excluded from the reported census totals.

\paragraph{Attribution and scope.}
The recovered target is the transitive double-omission branch
of Cameron--Fernandes--Leemans--Mixer, \emph{Highest rank of a
polytope for $A_n$}, arXiv:1605.09173v1, Proposition~3.4.
Its binary partition strategy is acknowledged in the separately
reviewed reduction. Here that reduction permits the stronger
local exclusion of all transitive five-mark subgroups of $W$,
without extending them to eight generators. Orbit enumeration,
the intersection property, and the recursive end criterion are
established tools; the latter is acquired from our reviewed F004
enumeration component. All subgroup calculations use GAP.
This packet proves only the displayed conditional branch. The
case where every double omission is intransitive, and the other
branches supplying the binary hypothesis, remain separate.
No whole-manuscript recovery, open-conjecture resolution, Lean
verification, or numerical compression factor is claimed here.
\end{document}
